\newcommand{\CurrentDetectorTable}{%
\begin{table*}[!t]
\caption{Held-validation detector operating point at severity 0.2.}
\label{tab:detector}
\centering\footnotesize
\setlength{\tabcolsep}{3pt}
\begin{tabular}{llrrr}
\toprule
Family & Scope & F1 & Recall & False positive (\%) \\
\midrule
clean & none & n/a & n/a & 0.03 \\
intermittent & missing & n/a & n/a & 0.05 \\
outage & missing & n/a & n/a & 0.00 \\
block & missing & n/a & n/a & 0.03 \\
noise & finite & 0.42 & 0.27 & 0.20 \\
bias & finite & 0.47 & 0.31 & 0.13 \\
drift & finite & 0.38 & 0.24 & 0.12 \\
stuck & finite & 0.66 & 0.55 & 0.40 \\
mixed & finite & 0.51 & 0.36 & 0.14 \\
\bottomrule
\end{tabular}
\par\vspace{2pt}
\begin{minipage}{\textwidth}\scriptsize Thresholds are frozen deployment buffers at the 5\% per-channel/context intact false-positive ceiling. Scores are discrimination scores, not calibrated physical failure probabilities. Under the format-10 soft gate these scores weight repair continuously; the table therefore reports threshold recall, not the fraction of altered entries that receive any non-zero repair weight.\end{minipage}
\end{table*}
}
\newcommand{\CurrentRecoveryTable}{%
\begin{table*}[!t]
\caption{Post-freeze recovery and detection benchmark by fault family.}
\label{tab:recovery}
\centering\footnotesize
\setlength{\tabcolsep}{3pt}
\begin{tabular}{llrrrrrr}
\toprule
Family & Scope & F1 & Recall & Reconst. RMSE & Intact disp. & False repair (\%) & Cov. (\%) \\
\midrule
clean & none & n/a & n/a & n/a & 0.017 & 0.0 & 98.4 \\
intermittent & missing & n/a & n/a & 0.46 & 0.021 & 0.1 & 98.5 \\
outage & missing & n/a & n/a & 0.46 & 0.011 & 0.0 & 98.4 \\
block & missing & n/a & n/a & 0.45 & 0.019 & 0.0 & 98.6 \\
noise & finite & 0.42 & 0.27 & 0.22 & 0.019 & 0.2 & 98.4 \\
bias & finite & 0.47 & 0.31 & 0.26 & 0.022 & 0.1 & 98.5 \\
drift & finite & 0.38 & 0.24 & 0.19 & 0.020 & 0.1 & 98.7 \\
stuck & finite & 0.66 & 0.55 & 1.99 & 0.023 & 0.4 & 97.6 \\
mixed & finite & 0.51 & 0.36 & 0.47 & 0.015 & 0.1 & 98.5 \\
\bottomrule
\end{tabular}
\par\vspace{2pt}
\begin{minipage}{\textwidth}\scriptsize Detection, false repair and intact displacement use the held selection recovery audit at severity 0.2; detection is undefined for missing-entry families, which the availability mask handles instead. Reconstruction and displacement are original-linear sensor units. Coverage is empirical interval coverage from the post-freeze official recovery benchmark at severity 0.2, where available. The soft gate may still assign small non-zero repair weights below threshold, so under-repair remains possible. Nothing was tuned on these values.\end{minipage}
\end{table*}
}
\subsection{\textbf{Detector Operating Characteristics}}
\CurrentDetectorTable
Table~\ref{tab:detector} makes the detector limitation explicit. Clean false
positives are 0.0322\% of intact entries, but finite-fault recall
is only 0.24 for drift,
0.31 for bias, 0.27 for
noise, and 0.55 for stuck. Recall is low in every
finite family, so the deployed path misses many genuine deviations rather than
over-repairing healthy cells.
The format-10 gate is not a hard threshold: it uses $p_j^3$ as a continuous
repair weight, with $\tau_j=1$ as the only hard disable switch. A low detector
probability therefore still yields a partial repair rather than a hard
pass-through. This mitigates
but does not remove the under-repair failure mode, and the paper makes no
universal recovery claim.

\subsection{\textbf{Recovery and Calibration by Fault Family}}
\CurrentRecoveryTable
Table~\ref{tab:recovery} pairs the detector audit with the complete recovery
matrix. The worst intact displacement is 0.02521
original-linear units, and the worst single-window intact RMSE is
0.06907. Stuck faults are the easiest finite
family to flag (F1 0.66) but still the hardest to
reconstruct, consistent with a held value being ambiguous rather than out of
range. Clean official coverage is 98.4\%; under
severity-0.2 stuck and mixed faults it is
97.6\%
and 98.5\%, respectively.
All 41 declared contexts completed and passed their
intact-cell guards.
